% ----------------------------------------------------------
\begin{frame}{Restriction $\sigma$}
  The restriction operator (also called \emph{selection}) \textbf{filters
  rows} of a relation, keeping only those tuples that satisfy a given
  Boolean condition.  It is the relational algebra equivalent of the SQL
  \texttt{WHERE} clause.

  \begin{block}{Key Points}
    \begin{itemize}
      \item The \emph{schema} of the output is identical to the input;
            only the number of rows may decrease.
      \item Conditions can use comparison operators
            ($=$, $\neq$, $<$, $>$, $\leq$, $\geq$) and logical
            connectives $\wedge$ (AND) and $\vee$ (OR).
      \item Multiple restrictions can be chained or merged:
            \[
              \sigma_{C_1}(\sigma_{C_2}(R)) = \sigma_{C_1 \wedge C_2}(R)
            \]
      \item Syntax: $\sigma_{\text{condition}}(R)$
    \end{itemize}
  \end{block}

  \begin{examples}{Example}
    \[
      \sigma_{\text{Department} = \text{`D01'}}(\text{Employees})
    \]
    Returns all employee tuples where the Department attribute equals
    \texttt{D01}.  All other tuples are discarded; the column set is
    unchanged.
  \end{examples}
\end{frame}

% ----------------------------------------------------------
\begin{frame}{Projection $\pi$}
  The projection operator \textbf{selects specific columns} from a
  relation, discarding all other attributes.  It corresponds to naming
  individual columns in a SQL \texttt{SELECT} list (as opposed to
  \texttt{SELECT *}).

  \begin{block}{Key Points}
    \begin{itemize}
      \item The number of rows stays the same \emph{unless} removing
            columns causes duplicate tuples — duplicates are then
            eliminated (pure relational algebra).  In SQL you need
            \texttt{DISTINCT} explicitly.
      \item Syntax: $\pi_{\text{col}_1,\,\text{col}_2,\,\ldots}(R)$
    \end{itemize}
  \end{block}

  \begin{examples}{Examples}
    \begin{itemize}
      \item $\pi_{\text{User, ID, Department}}(\text{Employees})$ —
            return only those three columns from the Employees relation.
      \item Combining restriction and projection:
            \[
              \pi_{\text{User}}\!\left(\sigma_{\text{Dept}=\text{`D01'}}(\text{Employees})\right)
            \]
            First filter to department D01, then project to the User
            column only.  This is a typical two-step pattern.
    \end{itemize}
  \end{examples}
\end{frame}

% ----------------------------------------------------------
\begin{frame}{Rename $\rho$}
  The rename operator changes the \textbf{name of a relation or its
  attributes} without modifying the actual data.  It is a purely
  structural operation needed whenever two relations that are to be
  combined share attribute names that would otherwise clash.

  \begin{block}{Key Points}
    \begin{itemize}
      \item \textbf{Rename relation:}
            $\rho_{\text{NewName}}(R)$ — gives the whole relation a new
            name.
      \item \textbf{Rename attribute:}
            $\rho_{(\text{NewAttr}/\text{OldAttr})}(R)$ — renames one
            attribute; all data remain unchanged.
      \item Particularly important when self-joining a table (you need two
            differently named copies) or when joining tables that happen to
            share a column name that should \emph{not} be the join
            criterion.
    \end{itemize}
  \end{block}

  \begin{examples}{Examples}
    \begin{itemize}
      \item $\rho_{\text{Employees}}(\text{Empl})$ — rename the relation
            \texttt{Empl} to \texttt{Employees}.
      \item $\rho_{(\text{Username}/\text{User})}(\text{Employees})$ —
            rename the attribute \texttt{User} to \texttt{Username} inside
            the Employees relation.
    \end{itemize}
  \end{examples}
\end{frame}

% ----------------------------------------------------------
\begin{frame}{Division $\div$}
  Division is the most complex basic operator.  $R_1 \div R_2$ returns
  all values of the \emph{non-shared} attributes of $R_1$ for which the
  corresponding shared attributes cover \textbf{every} tuple in $R_2$.
  In plain language: ``find all $X$ that are associated with
  \emph{all} $Y$.''

  \begin{block}{Formal Definition}
    Let $A$ denote the attributes in $R_1$ that are \emph{not} in $R_2$.
    Then:
    \[
      R_1 \div R_2 \;=\; \pi_A(R_1)
        \;\setminus\;
        \pi_A\!\bigl((\pi_A(R_1) \times R_2) \setminus R_1\bigr)
    \]
    This derivation shows that division is not strictly primitive — it can
    be expressed using projection, Cartesian product, and difference.
  \end{block}

  \begin{examples}{Example}
    \textbf{Families $\div$ Children} — suppose the Children relation
    lists two children: \emph{Linda} and \emph{Haley}.  The result
    contains every parent tuple from Families who has \emph{both} Linda
    and Haley as children.  A parent with only one of the two is
    excluded.
  \end{examples}
\end{frame}

% ----------------------------------------------------------
\begin{frame}{Additional Unary Operators}
  Basic relational algebra covers the operators introduced so far, but
  three additional unary operators are commonly added in practice to make
  the language more expressive and to bridge the gap to real SQL features.

  \begin{block}{Duplicate Elimination $\delta$}
    Removes all duplicate tuples from a relation — equivalent to SQL
    \texttt{DISTINCT}.\\
    Example: $\delta\!\left(\pi_{\text{Department}}(\text{Employee})\right)$
    — return a list of distinct department names.
  \end{block}

  \begin{block}{Grouping \& Aggregation $\gamma$}
    Groups tuples by one or more attributes and computes aggregate
    functions (\texttt{COUNT}, \texttt{SUM}, \texttt{AVG}, \texttt{MAX},
    \texttt{MIN}) over each group — equivalent to SQL \texttt{GROUP BY}.\\
    Example: $\gamma_{\text{Department,\;SUM(Salary)}}(\text{Employee})$
  \end{block}

  \begin{block}{Sorting $\tau$}
    Sorts the tuples of a relation by one or more attributes with optional
    \texttt{ASC}/\texttt{DESC} direction — equivalent to SQL
    \texttt{ORDER BY}.\\
    Example: $\tau_{\text{Department,\;Salary DESC}}(\text{Employee})$
  \end{block}
\end{frame}
